% backmatter/appendixD.tex —— 附录 D 可复现环境、版本与运行格式
% 数字身份纪律：E01单回合读数绑定E01-20260923T152230Z-90acda，
% E05v2索引读数绑定E05v2-20260929T143245Z-d7a8db；其余按所述归档与构造示例分记。
% 本附录描述的是本书实验台自身的落盘格式，不参与任何命题判决。
\chapter{可复现环境、版本与运行格式}
\label{app:repro}

本书的任何一条实测数字都要能被追回到一次具体运行。
本附录给出这条链路上每一环的格式：运行叫什么、落在哪里、每个字段是什么意思、
以及怎样先核对归档，再区分离线复算与新的模型运行。

\section{运行标识与目录}

\begin{sloppypar}
历史E01/E02/E03运行的标识由实验名、UTC时间戳和6位十六进制随机串拼成，
例如 \texttt{E01-20260923T152230Z-90acda}。
时间戳取自协调世界时，随机后缀用于降低同秒内运行重名的概率，不表示实验条件，也不保证绝不重名。
新运行还应在创建目录时拒绝已有路径；当前E11 v2与E12 v2入口使用\texttt{exist\_ok=False}，而历史工具并非全部如此。运行号生成与拒绝覆盖是两道不同的检查。
\end{sloppypar}

\begin{sloppypar}
历史E01/E02/E03运行在 \texttt{experiments/runs/<run\_id>/} 下主要保存三类工件：
\texttt{config.json}（这次运行的主要入参）、
E01/E02的训练轨迹位于\path{traces/<arm>/train/<task_id>-e<episode>/trace.json}，评价轨迹位于\path{traces/<arm>/sealed/<task_id>-r<repeat>/trace.json}。
E03的评价轨迹位于\path{traces/<split>/<arm>/<task_id>-r<repeat>/trace.json}，来源轨迹位于\path{traces/source/<split>/<task_id>-e<episode>/trace.json}（每回合一条轨迹）。另有
\texttt{summary.json}（聚合结果与判决）。
路径中的\texttt{arm}、\texttt{split}与\texttt{task\_id}分别表示实验臂、拆分与题号，\texttt{episode}和\texttt{repeat}分别表示训练回合序号与评价重复序号。
\texttt{experiments/results/}保留部分汇总及渲染副本，不能当作完整输入目录。从本版的\nolinkurl{analysis/inputs.json}查所用汇总的路径和哈希。公开预览的E03汇总副本替换了两处回溯中的临时目录，原件与副本分别登记哈希；核对数值可读副本，复核原始记忆的字符长度须读原件。拆分A的计数冲突见第~\ref{ch05}章。
\end{sloppypar}

\clearpage
\section{config.json：这次运行的入参}

\begin{longtable}{>{\raggedright\arraybackslash}p{3.6cm}>{\raggedright\arraybackslash}p{3.0cm}>{\raggedright\arraybackslash}p{6.6cm}}
\caption{运行入参：哪些字段决定了结果能不能复现}\label{tab:appd-config}\\
\toprule
字段 & 例 & 为什么必须记 \\
\midrule\endfirsthead
\multicolumn{3}{l}{表\thetable（续）}\\
\toprule
字段 & 例 & 为什么必须记 \\
\midrule\endhead
\midrule\multicolumn{3}{r}{续下页}\\\endfoot
\bottomrule\endlastfoot
\texttt{git\_head} & \texttt{89cec63...} & 代码版本入口；还需工作区差异与实际文件哈希 \\
\texttt{model} / \texttt{base\_url} & \texttt{deepseek-flash} & 模型与后端身份；别名相同不保证服务或权重版本相同 \\
\texttt{temperature} & \texttt{0.7} & 采样参数；未作为干预因素时应按协议保持一致 \\
\texttt{episode\_budget} & 调用 8 / 步数 12 & 任务内配额；等总预算还需加入任务外费用 \\
\texttt{E\_train\_episodes} / \texttt{R\_repeats} & \texttt{8} / \texttt{5} & 训练配额与每题重复次数；独立性须核对重置和依赖 \\
\texttt{base\_seed} & \texttt{20260924} & 记录标签的起点；不保证服务方使用同一种子 \\
\texttt{mde} / \texttt{unit} & \texttt{0.03} / \texttt{task} & 历史目标字段与重采样单位；须核对实际分析 \\
\end{longtable}

\section{trace.json：一个 episode 的回合记录}

下表说明历史静态执行器保留的轨迹字段。它记录模型文本、部分工具事件和最终判定，不是供应商完整响应或宿主全部副作用的逐项记录。

\begin{longtable}{>{\raggedright\arraybackslash}p{4.2cm}>{\raggedright\arraybackslash}p{9.4cm}}
\caption{单条轨迹的字段：谁跑的、跑了什么、花了多少、判定从哪来}\label{tab:appd-trace}\\
\toprule
字段 & 含义 \\
\midrule\endfirsthead
\multicolumn{2}{l}{表\thetable（续）}\\
\toprule
字段 & 含义 \\
\midrule\endhead
\midrule\multicolumn{2}{r}{续下页}\\\endfoot
\bottomrule\endlastfoot
\texttt{run\_id} / \texttt{agent} / \texttt{task\_id} & 回合标识、调用者名称与题号；父实验、臂与阶段仍须按目录和配置连接 \\
\texttt{seed} & 调用者传入的回合标签；生成规则随运行器与阶段而异，不能一律由\texttt{base\_seed}推回 \\
\texttt{started\_at} & 开始时间（UTC） \\
\texttt{steps} & 事件序列，以\texttt{kind}区分system / model / tool\_result / error / verdict \\
\texttt{usage} & 调用数、步数、提示与生成 token、墙钟秒数；\texttt{usd} 一律为 \texttt{null} \\
\texttt{budget} & 五项预算上限；停止须联查事件和终局判定，没有统一停止原因字段；token和墙钟阈值不是请求硬截断 \\
\texttt{verdict} / \texttt{verdict\_source} & 判定结果与来源；此处历史静态执行器调用 \texttt{task.verify}，新版合同须另核对 \\
\texttt{persistent\_writes} & 申报的版本层写入；旧静态执行器对非空列表抛异常，空列表不能证明不存在未登记写入 \\
\end{longtable}

\begin{sloppypar}
举一个真实的例子。E01 运行 \path{E01-20260923T152230Z-90acda} 中
\path{traces/grid4/sealed/classstate-07-r0/trace.json} 记录的那次 episode：
模型调用 \measured{3} 次、提示 token \measured{958}、生成 token \measured{84}、
墙钟\measured{2.14}秒，保存\measured{8}条事件记录；\texttt{usage.steps}为\measured{3}，对应执行循环次数，不等于事件条数。该回合判定为未通过，
判定来源写作 \texttt{task.verify（真实执行验收命令）}。
这几个数来自该回合归档。其\texttt{run\_id}为\path{static-20260923T154931Z-26925f}，不同于父实验的E01编号；\texttt{agent}为\texttt{grid4-sealed}。先由目录确定父实验、臂与阶段，再用配置和事件核对，不能只将两层\texttt{run\_id}作等值连接。\texttt{wall\_seconds}在最终\texttt{task.verify}之前取值，因此这段计时不包含最后一次验收的完整耗时。
\end{sloppypar}

\texttt{usage} 里的 \texttt{usd} 恒为 \texttt{null} 且 \texttt{unpriced} 为真：
后端不返回价格，本书的成本账只报实物量（调用数、token 数、墙钟时间），
折算美元还需能对应调用日期、模型和计价单位的价格记录。

旧后端在响应缺少用量字段时会将对应token默认记为零；没有保留原始响应时，不能据一个零字段认定真实用量已知为零。该静态轨迹也没有逐项保存外部任务题面、供应商完整响应或不可变模型修订身份。复核时应连接任务文件与配置，并把无法恢复的服务和用量信息保留为未知。

\section{复现步骤}

\begin{sloppypar}
复现首先分清目标：重算已有记录、使用脚本后端检查装置、还是重新发送付费模型请求。前两类不需要重新生成真实模型轨迹，后一类要检查数据范围、后端身份、额度和授权。先核对目标协议与实际文件，保留原运行；新版入口使用独立目录。不要在原始记录所在目录尝试“清理后重跑”，也不要把任务集重新生成当成读取旧记录的必要步骤。

\texttt{check\_backend.py}默认检查本地配置。只有显式指定\texttt{--probe}才会尝试一次真实请求，这次请求也应计入调用与费用。事后探测只能说明探测时的服务状态，核查历史运行须回到当时记录的模型与请求。

记录Git提交号有助于定位实现，但未提交改动、任务内容与协议也需单独保存。两个提交产生不同结果不自动排除复现失败；反过来，相同提交也不能保证远程模型、依赖或数据完全相同。应逐项解释差异，不用提交号单独判定可复现性。
\end{sloppypar}

\section{先离线复核，再决定是否重新运行}
\label{sec:appd-offline}

公开源码预览的离线入口列在\path{analysis/README.md}。下文以“完整档案”标出仅随维护仓库保存的命令；其余命令可在本次修订的公开预览中运行。复算E05案例和统计例题可以直接读取包内材料。若要查看E11或E12的默认预检报告，先在一份新解压的预览副本根目录运行：
\begin{verbatim}
python -m experiments.lab.build_taskset --no-lab-run
\end{verbatim}
这一步在\path{experiments/tasks/v1/}写入题库；已有同名文件可能被覆盖。本版默认只构建题库，显式使用\texttt{--no-lab-run}得到相同结果。只查历史汇总时无需生成题库。

\subsection{第一层：检查文件是否属于所述运行}

复核一个表格时，最小输入不是那张表格，而是运行标识、配置、原始记录与生成表格的规则。先确认引用路径确实存在，再确认文件内容与证据索引中的哈希一致。哈希一致只能说明字节未变，不证明这些字节真实、充分或没有遗漏；但没有这一层，后来发现的差异可能只是两个分析者读取了不同版本。

同一运行可以同时有原始轨迹、人工可读摘要与后续审计报告。它们的职责不同。原始轨迹保存当时观察到的事件，摘要按某套规则压缩事件，审计报告解释摘要可能遗漏或混淆的部分。当三者不一致时，应沿着字段定义寻找原因，不能简单选择数值最符合预期的一份。原始记录也可能缺失，遇到这种情形应登记缺口，而不是凭摘要重造一条看似完整的轨迹。

文件路径与语义标识也应分开。目录中一个序列可能按实验臂、复用格点和重复编号分层；单个任务标识只说明题目，不说明它在实验中的角色。第15章的两条相同题号轨迹可能分别处在来源阶段与评价阶段，输入记忆和预算也不同。使用题号连接表格前，至少还要核对运行、臂、阶段和重复标识，避免将不同回合错误配对。

对跨日期续跑，新增文件属于原运行的哪个部分，应有明确的父运行和配置对应。若后续解释指出某项模型版本信息不可恢复，就保留“不确定”，不从相同别名推导权重相同。能够成功加载配置，只证明配置语法可读；它不证明旧运行的服务条件已经重建。

\subsection{第二层：重算已记录的事件与数量}

本版提供一个只读取E05 v2归档的案例索引脚本。它检查六个完整序列，将每个事件连接到保存的轨迹，按序列核对调用数，再汇总主回合、独特题目与通过数。它不发送模型请求，不修改归档轨迹，也不把新版协议应用到旧记录上。运行后更新的是派生案例索引。

\begin{verbatim}
python calculations/reproduce_ch15_case.py -o reader-work/ch15.json
\end{verbatim}

\begin{sloppypar}
示例将索引另存为\path{reader-work/ch15.json}，目标已存在时拒绝覆盖；省略\texttt{-o}选项时沿用\path{analysis/ch15-case-index-reproduction.json}默认位置并更新该派生索引。归档复算应得到\measured{6}个序列、\measured{9}个主回合、\measured{48}次调用和\measured{7}次通过，并区分\measured{2}道独特题与\measured{1}个任务族。这些分母没有互相替代关系：九条轨迹不是九道独立题，七次通过也不是七次有效技能更新。
\end{sloppypar}

若脚本发现调用数与序列账不一致，应保存差异，并进一步查看哪些事件被纳入了两边。不要直接将一侧改成另一侧以让检查变绿。可能的原因包括字段原本是累计量、重试单列、附加修复被另记，或者记录确实不完整。查清字段语义后，才决定修订哪一份派生报告。错误消息是调查入口，不能只被当作需要消除的障碍。

通过数也要追溯到实际验收字段。回复中声称成功、工具动作没有报错、回合内一次测试成功与最终判定通过，分别指向不同事件。读者可以利用索引定位原始轨迹，再按事件顺序检查。索引本身没有推断模型为何行动，也没有证明某条被检索经验对最终答案具有因果作用。

\subsection{第三层：复算教学算术，不冒充模型重复实验}

附录B、D、F的平台记录、事件分母与历史成本算术，以及附录C的两动作数学例，可分别复算：
\begin{verbatim}
python calculations/reproduce_platform_appendices.py
python calculations/reproduce_rl_arithmetic.py
\end{verbatim}
两者另存本版分析文件。平台入口中的接口演示使用显式合成夹具，没有执行候选或请求模型；数学入口只导入旧脚本的纯计算函数，不覆盖旧的有日期报告。

完整档案保留了第4章反馈审计所需的逐回合轨迹与复算脚本：
\begin{verbatim}
python calculations/reproduce_feedback_audit.py
\end{verbatim}
它核对E02轨迹与汇总、正文配对差和栈案例，并比较错误方法与局部候选的容器状态。输出位于\path{analysis/feedback-audit-reproduction.json}。局部候选执行正确不代表Agent曾经执行该修改；未归档的中间评审原文仍然缺失。

完整档案中的以下脚本复核第5章的检索与来源成本：
\begin{verbatim}
python calculations/reproduce_memory_audit.py
\end{verbatim}
该入口读取SHA绑定的经验文件和任务规格，重建选中条目，再与归档提示逐项比较；它不调用模型，也不补造缺失的渲染用量。

第6章的操作接口与局部算术可单独复核：
\begin{verbatim}
python calculations/reproduce_memory_operations.py
\end{verbatim}
此入口只重放教学对象，结果另存\path{analysis/memory-operations-reproduction.json}，不覆盖旧检查日志。

第7章的四题选择偏差算例与来源表格算术可用以下入口核对：
\begin{verbatim}
python calculations/reproduce_context_audit.py
\end{verbatim}
它以精确分数计算等能力候选的最高开发分期望，另列ACE正文表4与附录表12的不同成本范围，以及GEPA最佳候选出现时点与总预算的比值。结果另存以下文件，不构成E04或外部论文实验复现：
\begin{quote}\ttfamily
analysis/\allowbreak context-audit-reproduction.json
\end{quote}

第8章的代码参考注入、来源拒绝与费用图使用以下入口：
\begin{verbatim}
python calculations/reproduce_skill_audit.py
\end{verbatim}
它只读随书E05 v2轨迹，输出\path{analysis/skill-audit-reproduction.json}，并复核教学成本数组。报告按原运行配置逐项校验14份纯文本快照的字节，包括当时的源码和协议。任一摘要不符，复算即停止。旧E05累计账未纳入冻结包，历史分析脚本不属于本版可直接复算入口，依赖它的数值表已撤下。

\begin{samepage}
第9章的来源表格差值、独立候选覆盖率、排程反例与摊销算例可复算：
\begin{verbatim}
python calculations/reproduce_workflow_arithmetic.py
\end{verbatim}
\end{samepage}
输出位于\path{analysis/}\allowbreak\path{workflow-arithmetic-reproduction.json}。这些计算不产生E06结果，也不代替来源论文的实际运行。

完整档案中，第10章的SEAL表格差值、教学梯度、成本交点及区间判据可用以下脚本复算：
\begin{verbatim}
python calculations/reproduce_parameter_arithmetic.py
\end{verbatim}
结果位于\path{analysis/}\allowbreak\path{parameter-arithmetic-reproduction.json}。它只做确定性计算，没有训练参数模型，也不产生E08能力证据。

完整档案中，第11章的课程漂移、过滤分母、F1与格式奖励合账及来源表格差值可复算：
\begin{verbatim}
python calculations/reproduce_curriculum_arithmetic.py
\end{verbatim}
输出位于\path{analysis/}\allowbreak\path{curriculum-arithmetic-reproduction.json}，保持教学设定与论文报告的身份，不产生E07或E09训练结果。

第13章的历史触发信号、暖机占位、精确概率、三策略落点与成本算例可复算：
\begin{verbatim}
python calculations/reproduce_timing_audit.py
\end{verbatim}
结果另存于\path{analysis/}\allowbreak\path{timing-audit-reproduction.json}。它核对固定历史输入和教学模型，不执行E11，也不能估计替代调度策略的闭环成绩。

第14章的噪声比较、精度规划、余弦例题与单文件恢复可从归档数据复核。该命令也核对历史权限观察的来源身份。
\begin{verbatim}
python calculations/reproduce_reliability_audit.py
\end{verbatim}
输出位于\path{analysis/}\allowbreak\path{reliability-audit-reproduction.json}。旧E12的记账观察附有源码节选和校验值，入口按数据读取这些材料。重新评价候选执行权限需要单独的隔离环境。

完整档案另有脚本复算第15章的E01范围、任务族差值、读取者成本比例、题级点估计反号和E05令牌差值：
\begin{verbatim}
python calculations/reproduce_case_arithmetic.py
\end{verbatim}
结果位于\path{analysis/}\allowbreak\path{case-arithmetic-reproduction.json}。它区分80个封存评价回合与88个计时主回合，不把既有不同运行拼成一次集成实验。

完整档案另存了第16章语料比例、振荡、条件风险、零故障上限与独立选择的复算脚本：
\begin{verbatim}
python calculations/reproduce_scope_arithmetic.py
\end{verbatim}
结果保存为：
\begin{quote}\ttfamily
analysis/\allowbreak scope-arithmetic-reproduction.json
\end{quote}
它还以构造输入检查FGT与BWT使用不同参照，不产生本书的长期保持率或实测故障概率。

第10—16章同时包含来源论文报告值、真实运行记录和人为教学例。教学例用于说明比较口径，必须可以计算，也必须保持其人为设定的身份。完整档案的下列脚本复算成本直线交点、难度代理、混合分布均值、父代权重、效用、折扣后代增长、交互差和长期概率例。

\begin{verbatim}
python calculations/reproduce_teaching_arithmetic.py
\end{verbatim}

输出中的“通过”表示公式结果与保存的教学数值相符。它不表示E07、E08、E10或长期运行已经成功。脚本中的成本单位、条件故障概率与独立性都来自各例明确给定的条件；更改条件会改变答案，不能拿计算结果为条件本身作证。例如，独立比较的错误概率例只能解释选择风险，不能用来估计本书各项实际比较的误判概率。

复算也不代替公式来源核对。父代选择权重即使计算无误，若使用了不属于该算法版本的资格条件，仍然是在精确计算错误问题。第12章因此分别登记DGM与HyperAgents的合格集合、缩放中点和子代计数。公式身份、输入身份和算术结果需要同时正确，任何一项缺失都不能靠另外两项补齐。

对于有舍入的报告，应优先使用保存的未舍入量计算，最后再按表格精度显示。若手工用已经四舍五入的标准差反推样本数，结果可能与原计算相差一个整数。应该解释差异来自哪一层舍入，而不把它当成实验结果变化。附录A的样本量核算与本附录的教学核算，正是两类不同来源的派生量。

\subsection{第四层：检查展示文件来自当前源稿}

正文正确而PDF过期，是交付链中另一类可复现性问题。本版构建入口把当前书稿与依赖复制到临时构建目录，生成结果后留下诊断文件。检查时应读取诊断指定的实际日志，而不是仅凭主目录中某个旧日志的时间。编译退出成功也要同时核对交叉引用、缺字和图表资源是否来自当前源文件。过高浮动表格可能被裁掉而没有编译错误；除\texttt{Overfull}以外，还应检查过高浮动体警告。完整档案附有以下构建审计脚本。

\begin{verbatim}
python calculations/audit_book_build.py
\end{verbatim}

这个脚本需要用于PDF检查的依赖，核对当前PDF、实际构建快照与书稿源文件的对应关系，并检查构建状态。它与归档复算、教学算术入口使用不同的工具链。若正文刚改而新构建尚未结束，快照比较失败是合理的提醒：当前展示物仍对应旧稿。应等待真实构建完成再核对，而不是关闭该比较。

版面阅读仍不可省略。没有文字溢出，不意味着图中箭头指向正确；引用编号存在，不意味着来源真的支持那句话；数学符号能够排版，也不意味着单位一致。自动检查负责发现一部分可形式化问题，页面抽检负责发现视觉问题，来源与推理核对负责判断内容。三者的通过范围应分别报告。

\subsection{什么结果需要新运行编号}

若只是用相同规则重新汇总原始记录，新产物是旧运行的派生分析。若更换分析口径，例如重采样层级或缺失处理，应登记新的分析版本，而不制造新的模型运行编号。若重新向模型发出请求，即使题目和提示相同，也产生新的执行事件，需要保存新轨迹及其与旧设计的关系。这三种版本变化不应混用。

修复运行器后更需要新旧分离。E05 v3修改了格式处理、预算生命周期和技能准入行为；它的未来运行将回答新版装置的问题，不能倒替旧版中失败的几个回合。将成功旧回合与修复后的新回合拼成一套“完整新版”，会同时隐藏修复前看到过哪些结果，以及重新采样付出了多少资源。

若目的是复现外部方法，读者还需准备相应模型、环境、数据与预算。本书的短程代码实验台不能仅通过更改方法名称就变成参数训练、长期研究或完整软件工程基准。可以先做一个受限原型，但应将其标成原型并列出删去的机制。教学实现便于理解，不自动继承原论文全部结论。


\section{运行E01新版的离线检查}

以下一组命令需要完整档案。从仓库根目录启动，可以复查附录B的控制流与文件边界，不需要模型密钥。当前版本使用标准库构造输入；默认入口只打印就绪报告，显式选择离线演示后才创建新的本地记录目录。它没有付费执行选项，也不能用一个参数把构造结果切换成真实能力实验。

\begin{verbatim}
python -m unittest experiments.lab.test_e01_v2_contracts
python -m unittest experiments.lab.test_e01_v2_offline
python -m unittest experiments.lab.test_e01_v2_file_boundary
python -m unittest experiments.lab.test_e01_v2_analysis
python -m unittest experiments.lab.test_e01_v2_manifest
python -m unittest experiments.lab.test_e01_v2_execution
python -m unittest experiments.lab.test_e01_v2_evaluation_journal
python -m unittest experiments.lab.test_e05_v4_tools
python -m experiments.lab.run_e01_v2
python -m experiments.lab.run_e01_v2 --offline-demo
python calculations/demo_e01_final_journal.py
\end{verbatim}

前七组分别检查E01设计合同、调度与整条流程、公开文件边界、描述分析、任务清单、外部验收输入输出及持久评测日志；第八组检查E05 v4草案的独立文件工具。它们都不证明候选进程隔离已经完成；失败时应先保留报错与版本，不把旧输出复制过来充当本次通过。演示每次使用新的编号，默认落在本仓库的\nolinkurl{experiments/offline/}下，目录已存在时拒绝覆盖。读者可以通过入口的帮助查看输出目录选项，将教学记录保存在自己指定的位置。

最后一条命令单独演示持久日志的预留、重开、清理对账与缺失保留，每次另建目录；它与前一条序列记忆演示的记录布局不同。两种演示都不执行候选代码。结果顶层明确写有合成工程记录身份、真实API调用为零以及尚未实现的边界。序列记忆演示的各任务目录分别保存提交、构造评价、资源事件和汇总，便于检查提交是否先于评价、一次逻辑生成是否包含多次请求尝试。不要把演示目录放进真实实验结果汇总脚本，也不要把预设失败改成成功后仍声称复现了同一份流程。

\begin{sloppypar}
\begin{samepage}
两版均收录以下命令，可核对旧文件组件的历史观察及其来源。
\begin{quote}\ttfamily\small
calculations/\allowbreak reproduce\_reliability\_audit.py
\end{quote}
\end{samepage}
权限观察来自当时保存的临时标记检查报告。入口核对其字节身份；当前文件工具的回归测试另行验证目录访问规则。进程能够读取哪些文件，还取决于运行它的账户和隔离环境。
\end{sloppypar}

\subsection{公开提炼与近期读取的独立演示}

第一条是完整档案的单元测试，第二条是两版都收录的演示：
\begin{verbatim}
python -m unittest experiments.lab.test_e01_v2_public_memory
python calculations/demo_e01_public_memory.py
\end{verbatim}

本版的公开记忆、准入凭据与持久准入演示均使用下列文件中的新接口夹具：
\begin{quote}\ttfamily\small
calculations/\allowbreak platform\_demo\_fixtures.py
\end{quote}
这些接口夹具与仓库保留的历史JSON任务目录是不同输入；不能将夹具冒充旧任务或其验收答案。演示检查的是事件顺序与记录关联，不应为运行演示而覆写历史任务目录。

第二条只向标准输出打印构造请求、来源记录及读取快照，不保存正式实验结果，不发送模型请求，也不执行候选。任务标签、公开请求、回复与笔记均为手写的接口夹具。应检查公开提炼载荷不含宿主文件，在线读取逐次替换，迁移快照保持固定；不能把这些结构观察当作模型解决任务或生成了合格经验。旧演示与本演示采用不同读取配置，分别保留身份。

\subsection{复核来源、回复与准入记录的关联}

第二条演示和第四条数学复核可在公开预览运行；第一条单元测试和第三条历史触发审计需要完整档案。
\begin{verbatim}
python -m unittest experiments.lab.test_e01_v2_admission_receipts
python calculations/demo_e01_admission_receipts.py
python calculations/audit_ch13_trigger_replay.py
python calculations/audit_ch14_cosine_bound.py
\end{verbatim}

前两条不发送请求或执行候选。读者可以检查候选是否等于保存回复中指定的片段、审核引用的摘要是否一致、发布后快照是否包含同一条目。示例的审核类别是脚本夹具，不是实际语义审核；不能将其改标为模型结果或人类批准。该组件尚无磁盘重开协议，也没有接入真实运行器。

第三条只读旧轨迹并复算第\ref{ch13}章的触发信号，区分暖机占位值与有效窗口；它不会重跑任务或覆盖旧结果。第四条以精确分数复核第\ref{ch14}章的余弦差反例，检验的是一个数学推论所需的条件，不是模型性能或原文系统的复现。两条审计均只向标准输出返回结果，正式复核时应另存带时间的记录，并保留旧文件身份。


\begin{samepage}
\subsection{复核一次经验准入在中断后怎样恢复}

前节的内存组件保持原样；新增持久日志将一次提取的派发机会、回复、用量、候选片段、审核和结果快照联系起来。以下两条命令需要完整档案，均使用构造数据，不发送模型请求：
\begin{verbatim}
python -m unittest experiments.lab.test_e01_v2_admission_journal
python calculations/demo_e01_admission_journal.py
\end{verbatim}
\end{samepage}

演示先保存调用与派发预留，关闭数据库后重新打开，确认未完成请求阻止重发。随后填入明确标作脚本夹具的回复与用量，再经一次重开形成候选和审核，最后保存记忆快照。这里的关闭和重开是恢复测试，不是声称服务商恰好执行一次；真实请求可能在本地未拿到回复时已经完成。遇到这种不确定状态，程序保留未知并等待协调，不能用一次新成功响应覆盖它。

输出中的token数是人为检查值，不能计入本书实测成本。数据库事务保证的是本组件所保存的关联一致性：成功回复与该次用量一同提交，旧句柄不能重复结算，重复发布不能增加条目。审核仍是宿主提供的声明，既不证明经验正确，也不认证审核者身份。返回Memory快照也不等于已恢复部署进程、工具缓存和外部效果。

这一单元只限制一次online提取的资源，不负责跨任务的总预算。为下一任务新建数据库，不能被当作取得一份免费配额。独立的离线编排组件已连接运行身份、共享预算权威、持久序列游标与合成准备记忆；它只将U臂提取纳入共享池，没有覆盖准备、求解、测试与重试的完整费用。真实准备记忆、实际候选执行和服务适配仍需接通。

\subsection{序列编排的离线检查}

完整档案保留了三组检查，分别覆盖跨单元身份与预算游标、M0及记忆合并、组件编排：
\begin{verbatim}
python -m unittest experiments.lab.test_e01_v2_orchestration
python -m unittest experiments.lab.test_e01_v2_m0_memory
python -m unittest experiments.lab.test_e01_v2_assembly
\end{verbatim}
检查使用构造材料和本地临时记录，不发送模型请求。完成单元在重开后不再支取预算、不重复合并记忆，是局部恢复不变量；它不保证任意中途退出后能够自动续跑，也不证明整个实验的真实token与墙钟已被硬封顶。单次准入的发布快照与多次累计后的存储须按各自定义核对，不能仅因某条笔记进入存储就称整个快照相等。

\section{历史运行环境与本版构建条件}

\begin{itemize}
  \item 历史环境日志记载Python\texttt{3.13}（见完整档案\path{agent/nightly/2026-09-24.md}）；这不是每条轨迹各自记录的解释器身份。实验运行、论文抽取和书稿审查的依赖不同，应分别记录。
  \item 历史验收命令使用标准库\texttt{unittest}。更换测试框架前应核对发现规则、失败处理与退出码，再判断新旧结果是否可合并。
  \item 模型后端经OpenAI兼容接口调用，本书运行入口从环境变量或\texttt{experiments/lab/.env}读取凭据；凭据不作为配置和轨迹的主动记录字段。
  \item \begin{sloppypar}
        后端对指定网络异常及HTTP状态重试，例如读超时、429、500、502、503和504。本包保留的运行器可以把\texttt{backend\_retries\_total}写入配置，但归档E01/E02配置及其汇总配置没有该字段，不能从当前源码补称当时已记录。E03续跑配置中的计数也只反映对应后端实例，不能覆盖早期进程。历史\texttt{usage.calls}在逻辑生成成功返回后增加一次，内部HTTP尝试不逐次进入该字段；没有完整请求账时，重试总数和失败用量应保持缺失。
        \end{sloppypar}
  \item 本版经 \texttt{book/build.ps1} 调用TeX Live 2023的XeLaTeX；该入口采用临时构建目录绕开已遇到的中文路径问题，并保存实际构建诊断。
\end{itemize}
